\documentclass[10pt,a4paper]{amsart}
\usepackage[margin=19mm]{geometry}
\usepackage{amsmath,amssymb,mathtools}
\usepackage[scr=rsfs,cal=euler]{mathalfa}
\usepackage{xfrac}
\usepackage[unicode,hidelinks,bookmarks=false]{hyperref}
\newcommand{\ie}{i.e.\ }
\newcommand{\cf}{cf.\ }
\newcommand{\resp}{respectively\ }
\newcommand{\Subsection}{\subsection}
\providecommand{\nobreakdash}{\nobreak\hbox{-}}
\setlength{\emergencystretch}{2em}
\multlinegap0pt
\usepackage{bm}
\usepackage{bbm}
\usepackage{dsfont}
\usepackage{xspace}
\usepackage{cancel}
\usepackage{mathtools}
\usepackage{amsthm}
\makeatletter
\@ifundefined{thm}{\newtheorem{thm}{Thm}}{}
\@ifundefined{lemma}{\newtheorem{lemma}[thm]{Lemma}}{}
\makeatother
\newcommand{\E}{\mathbb{E}}
\title[Robust Learning of a Group DRO Neuron]{Robust Learning of a Group DRO Neuron}
\author{Guyang Cao, Shuyao Li, Sushrut Karmalkar, Jelena Diakonikolas}
\date{Source: arXiv:2601.18115v2 (CC BY 4.0)}
\begin{document}
\maketitle

\noindent\emph{Source and licence.} Robust Learning of a Group DRO Neuron. Version arXiv:2601.18115v2 (CC BY 4.0), distributed under the Creative Commons Attribution 4.0 International licence, as recorded in the arXiv licence metadata for this exact version and shown on the abstract page https://arxiv.org/abs/2601.18115v2. Full rights notice: licenses/arxiv-2601.18115-CC-BY-4.0.txt.\par\smallskip

\noindent\textit{Editorial scope.} Counted unit: the Lemma whose source label is \texttt{lemma:concentration-via-clipping} in the article (source0.tex lines 1054--1067, source label \texttt{lemma:concentration-via-clipping}), delivered together with the article's own complete original proof of it (source0.tex lines 1068--1180, one \texttt{proof} environment of 113 lines). No mathematics is written, repaired or re-derived: statement and proof are the article's own text line for line. Required context retained uncounted: none. Every definition, display and label of those lines is retained unchanged. Omitted from this package and not redistributed: 1--1053, 1181--2046. Nothing inside a retained excerpt is omitted or altered: every retained body is delivered exactly as the article's lines are written, so no omission receipt of any kind accompanies this package. Citations: the retained excerpts carry no citation command at all, so no bibliography entry of the article is transcribed and no reference is redistributed. Reference adaptation: none was needed and none was made; every reference inside the delivered unit resolves inside it, so no unresolved reference or citation is produced. Printed slips kept literal: no printed word, symbol, hypothesis or number of the counted statement or of its proof was altered. Layout and class: the article's own class is \texttt{article} with packages including geometry, siunitx, multirow, natbib, listings, xcolor, caption, algorithm2e, inputenc, fontenc, url, booktabs, amsfonts, nicefrac; the wrapper is the lane's pinned \texttt{amsart} editorial wrapper at 10pt on A4, which restates the theorem-like environments the retained text needs and adds the article's own macro definitions that the retained text needs (\texttt{\textbackslash E}), with extra wrapper packages (bm, bbm, dsfont, xspace, cancel, mathtools). The delivered numbering is the unit's own. Assets: no retained span contains a figure environment, a graphics-inclusion command, a PDF-page inclusion, a table or a TikZ picture; no image file is copied into this package, the package declares no asset dependency and no image file is redistributed with it. Licence: version arXiv:2601.18115v2 of this article is distributed under the Creative Commons Attribution 4.0 International licence, as recorded in the arXiv licence metadata for this exact version and shown on the abstract page https://arxiv.org/abs/2601.18115v2. A full rights notice is delivered at licenses/arxiv-2601.18115-CC-BY-4.0.txt. Characters: every retained excerpt is pure ASCII, so the delivered engine, XeLaTeX with its default Latin Modern text font, reports no missing character.
\begin{lemma}[One-Dimensional Concentration via Clipping]\label{lemma:concentration-via-clipping}
Let $Z$ be a real-valued random variable with mean $\mu = \E[Z]$. 
Assume $Z$ satisfies the tail bound $\Pr(|Z| > t) \le 2\exp(-t^{1/\tau}/B)$ for all $t \ge t_0$, where $t_0 > 0$ and $\tau \ge 1$ is a bounded integer. 
Let $Z_1, \dots, Z_n$ be $n$ i.i.d. copies of $Z$.

For any error $\epsilon > 0$, failure probability $\delta \in (0, 1)$, and sufficiently large $n$ satisfying  
\[
n = \tilde{O}\left( \max\left( \frac{B^{2\tau}}{\epsilon^2}, \frac{t_0^2}{\epsilon^2} \right) \log \left( \frac 1 \delta \right) \right),
\]
with probability at least $1-\delta$, we have:
\[
\left| \frac{1}{n}\sum_{i=1}^n Z_i - \mu \right| \le \epsilon.
\]
\end{lemma}

\begin{proof}
Let the clipping threshold be $S \geq t_0$, to be determined later, and define $Z_c := \max(-S, \min(Z, S))$. 
We decompose the total error using the triangle inequality:
\[
\left|\frac{1}{n}\sum_{i=1}^n Z_i - \mu\right| \le \underbrace{\left|\frac{1}{n}\sum_{i=1}^n (Z_i - Z_{i,c})\right|}_{\text{(I) Clipping Error}} + \underbrace{\left|\frac{1}{n}\sum_{i=1}^n Z_{i,c} - \E[Z_c]\right|}_{\text{(II) Concentration Error}} + \underbrace{|\E[Z_c] - \mu|}_{\text{(III) Bias Error}}.
\]
We now bound each of these three terms.

\textbf{Clipping term:} 
The first term is non-zero only if at least one sample $|Z_i| > S$. 
Let $\mathcal{E}_{\text{clip}}$ be this event. 
By a union bound over the samples and the tail bound on $Z$:
\[
\Pr(\mathcal{E}_{\text{clip}}) = \Pr(\exists i: |Z_i| > S) \le n \cdot \Pr(|Z| > S) \le 2n \exp(-S^{1/\tau}/B).
\]
This means with probability $1-2n \exp(-S^{1/\tau}/B)$, we have $\left|\frac{1}{n}\sum_{i=1}^n (Z_i - Z_{i,c})\right| = 0$.

\textbf{Bias term:} The bias $|\E[Z_c] - \mu| = |\E[Z_c - Z]|$ is bounded by the integral of the tail probability:
\[
|\E[Z_c - Z]| \le \E[|Z| \cdot \mathbb{I}(|Z|>S)] = \int_S^\infty \Pr(|Z|>t)dt \le \int_S^\infty 2e^{-t^{1/\tau}/B}dt,
\]
where the first inequality follows from the natural coupling between $Z_c$ and $Z$ assigning every point in $Z_c$ to the corresponding unclipped point in $Z$. 
If the point is unclipped, this contributes $0$,
if the point is clipped, then this means $|Z| > S$ and it contributes $|Z|-S < |Z|$. 
The equality is the standard way to express expectations of nonnegative random variables as tail integrals. 
The final inequality follows from the tail bound on $|Z|$. 

We now further upper bound the right hand side above.

Let $u = t^{1/\tau}/B$, so $t = (Bu)^\tau$ and $dt = \tau B (Bu)^{\tau-1} du$. 
The lower limit of integration becomes $u_0 = S^{1/\tau}/B$.
\[
\int_{u_0}^\infty 2e^{-u} \tau B (Bu)^{\tau-1} du = 2\tau B^\tau \int_{u_0}^\infty u^{\tau-1} e^{-u} du.
\]
For $u_0 \ge \tau-1$ (which will be satisfied by the choice of \(S\) later), 
we can bound $\int_{u_0}^\infty u^{\tau-1} e^{-u} du \le c_\tau u_0^{\tau-1} e^{-u_0}$ for some constant $c_\tau < 2$ depending on $\tau$ 
(by standard facts about exponential integrals). 
The bias is thus bounded by $2\tau B^\tau c_\tau u_0^{\tau -1} e^{-u_0} = 2\tau B^\tau c_\tau (S^{1/\tau}/B)^{\tau -1} e^{-(S^{1/\tau}/B)} = \Theta_\tau$.


\textbf{Concentration term:} Finally, we bound the concentration term. The clipped variable $Z_c$ is bounded in $[-S, S]$. We can apply Hoeffding's inequality to its empirical mean. With probability at least $1-\delta/2$:
\[
\left|\frac{1}{n}\sum_{i=1}^n Z_{i,c} - \E[Z_c]\right| \le \sqrt{\frac{2S^2\log(4/\delta)}{n}} = S\sqrt{\frac{2\log(4/\delta)}{n}}.
\]
Substituting $S = (B \log(4n/\delta))^\tau$, this deviation is bounded by $(B \log(4n/\delta))^\tau \sqrt{\frac{2\log(4/\delta)}{n}}$.

Accounting for everything, we see that, with high probability, the deviation is at most the sum of the concentration and bias errors. We need to choose the clipping threshold $S$ and the sample size $n$ to make this total error at most $\epsilon$, while keeping the failure probability at most $\delta$.

This requires satisfying the following conditions simultaneously: 
\begin{align}
    S &\ge (B \log(4n/\delta))^\tau \label{eq:cond_clip} \\
    c_\tau B S^{(\tau-1)/\tau} e^{-S^{1/\tau}/B} &\le \epsilon/2 \label{eq:cond_bias} \\
    n &\ge \frac{8 S^2}{\epsilon^2} \log(4/\delta) \label{eq:cond_conc}\\
	S &\ge t_0\\
    S &\ge (\tau-1)^{\tau} B^{\tau} \label{eq:u_0}
\end{align}
where the first condition ensures the clipping event probability is at most $\delta/2$, the second bounds the bias (where the constants have been absorbed into $c_\tau$), and the third bounds the concentration error of the clipped estimators. The second-last condition is just to ensure that the tail bound applies, and the final condition corresponds to $u_0 \geq \tau-1$. 

\textbf{Choosing $S$:} To satisfy all of these inequalities simultaneously, we define the clipping threshold $S$:
\[
S := \max\left( \left(B \log\left(\frac{C_\tau B^\tau}{\epsilon}\right)\right)^\tau, \left((\tau-1) ~B \log(4n/\delta)\right)^\tau, t_0 \right)
\]

for a sufficiently large constant $C_\tau$. The first term in the $\max$ is chosen to satisfy the bias condition \eqref{eq:cond_bias}, and the second term satisfies the clipping condition \eqref{eq:cond_clip} as well as \eqref{eq:u_0}. 
We now explain how to derive the first term.

Let us define the variable $X = S^{1/\tau}/B$. The inequality can be rewritten in terms of $X$ by substituting $S^{1/\tau} = BX$:
\begin{align*}
c_\tau B (BX)^{\tau-1} e^{-X} &\le \frac{\epsilon}{2} \\
c_\tau B^\tau X^{\tau-1} e^{-X} &\le \frac{\epsilon}{2}
\end{align*}
To satisfy this, we need to find a sufficiently large $X$. For any $\tau \ge 1$, there exists a (constant) threshold $X_0(\tau)$ such that for all $X \ge X_0(\tau)$, the polynomial term $X^{\tau-1}$ is bounded by an exponential, i.e., $X^{\tau-1} \le e^{X/2}$. Assuming $X$ is large enough to meet this condition, our inequality becomes:
\[
c_\tau B^\tau e^{X/2} e^{-X} = c_\tau B^\tau e^{-X/2} \le \frac{\epsilon}{2}
\]
Solving for $X$:
\begin{align*}
e^{-X/2} &\le \frac{\epsilon}{2 c_\tau B^\tau} \\
-X/2 &\le \log\left(\frac{\epsilon}{2 c_\tau B^\tau}\right) = -\log\left(\frac{2 c_\tau B^\tau}{\epsilon}\right) \\
X &\ge 2 \log\left(\frac{2 c_\tau B^\tau}{\epsilon}\right)
\end{align*}
This choice of $X$ is consistent with the initial assumption that $X$ is large (for small $\epsilon$), as the logarithm will be large. We can absorb the constant factors into a new constant $C_\tau$ and state that the condition is satisfied if
$
X \ge 2\log\left(C_\tau \frac{B^\tau}{\epsilon}\right).
$
Substituting back $X = S^{1/\tau}/B$, we get:
\[
\frac{S^{1/\tau}}{B} \ge 2\log\left(C_\tau \frac{B^\tau}{\epsilon}\right)
\]
This gives the required lower bound on $S$:
\[
S \ge \left(2B \log\left(C_\tau \frac{B^\tau}{\epsilon}\right)\right)^\tau
\]
This demonstrates that a choice of $S$ satisfying $S = \tilde{O}((B \log(B/\epsilon))^\tau)$ is sufficient to control the bias term.

\textbf{Solving for $n$:} Substituting this choice of $S$ into the concentration condition \eqref{eq:cond_conc} yields a lower bound on $n$. This gives three requirements on $n$, corresponding to which term in the $\max$ defining $S$ is dominant.
\begin{enumerate}
    \item If $S$ is determined by the bias term (the first term), then $n$ must satisfy:
    $n \ge \frac{8 \log(4/\delta)}{\epsilon^2} \left(B \log\left(\frac{C_\tau B^\tau}{\epsilon}\right)\right)^{2\tau}.$
    This contributes to $n = \tilde{O}\left(\frac{B^{2\tau}}{\epsilon^2} \log\left(\frac 1 \delta \right)\right)$ in the stated sample complexity.
    
\item If $S$ is determined by the clipping probability term (the second term), we get a recursive condition on $n$:
    \[ n \ge \frac{8 \log(4/\delta)}{\epsilon^2} (B \log(4n/\delta))^{2\tau}. \]
    This is a recursive inequality of the form $n \ge K \log^{2\tau}(n)$ where $K = \frac{8 B^{2\tau} \log(4/\delta)}{\epsilon^2}$; by the standard technique for solving such recursive inequalities this is satisfied by all $n \geq \max((2\tau e)^{2\tau}, 2K\log^{2\tau}(2K)) \allowbreak = \Theta(K\log^{2\tau}(K))$, and so it suffices that:
    \[ n = \tilde{O}\left(K \log^{2\tau}(K)\right) = \tilde{O}\left(\frac{B^{2\tau}}{\epsilon^2} \log^{2\tau}\left(\frac{B^{2\tau}}{\epsilon^2}\right)\right). \]
    Since the $\tilde{O}$ notation absorbs polylogarithmic factors, this simplifies to $n = \tilde{O}\left(\frac{B^{2\tau}}{\epsilon^2} \log(\frac 1 \delta)\right)$.
    
    \item If $S$ is determined by the threshold term $t_0$ (the third term), then $S = t_0$ and we need:
    $ n \ge \frac{8 \log(4/\delta)}{\epsilon^2} t_0^2. $
    This contributes $n = O \left(\frac{t_0^2}{\epsilon^2} \log(1/\delta)\right)$ to the sample complexity.
\end{enumerate}
The total sample complexity is the maximum of these three requirements, giving us the sample complexity stated in the lemma. 
\end{proof}

\end{document}
